\documentclass[12pt,a4paper]{article}
\usepackage[utf8]{inputenc}
\usepackage{geometry}
\geometry{a4paper, margin=1in}
\usepackage{hyperref}
\usepackage{titlesec}
\usepackage{graphicx}
\usepackage{booktabs}

\title{\textbf{EvoCracker: Phase 6 Technical Report \\ \large System Stabilization \& AI Pathfinding Refinement}}
\author{Development Team}
\date{\today}

\begin{document}

\maketitle

\begin{abstract}
This report outlines the comprehensive architectural improvements and bug resolutions implemented in the EvoCracker system during Phase 6. The primary focus of this phase was to address critical discrepancies in the AI Analytics Dashboard, enforce consistent enemy spawning, and resolve mathematical deadlocks within the core pathfinding algorithms (Depth-Limited Search, Iterative Deepening Search, and Hill Climbing). These resolutions stabilize the evolutionary pipeline and ensure absolute synchronization between the simulation and its telemetry.
\end{abstract}

\section{Introduction}
Phase 6 development was initiated to resolve severe inconsistencies identified during the integration of the Genetic Algorithm (GA) with the game simulation. The telemetry data presented in the AI Analytics Panel frequently decoupled from the rendering engine, presenting inaccurate metrics. Furthermore, specific enemy archetypes exhibited catastrophic movement failures, appearing "frozen" or locked in loops when transitioning into alert states. This document details the exact nature of these defects and the technical solutions applied to rectify them.

\section{Analytics \& Core Engine Resolutions}

\subsection{Analytics \& Sprite Synchronization}
\textbf{Issue:} The telemetry dashboard displayed "made-up" metrics and inaccurate enemy names relative to their visual sprites. The GA was optimizing entities that did not correspond to their visual representation. \\
\textbf{Resolution:} The \texttt{createEnemySprite} method within \texttt{SpriteFactory.ts} was refactored. The hardcoded fallback to generic dungeon-pack sprites was removed. It now dynamically maps the runtime \texttt{EnemyType} string directly to its corresponding \texttt{CharacterDef} index. The rendered sprite now strictly correlates with the exact name and data recorded in the analytics snapshot.

\subsection{Guaranteed Genetic Diversity (Spawning Logic)}
\textbf{Issue:} The procedural generator gated enemy types behind specific floor thresholds (e.g., Floor 1 only generated Toads and Ghosts). This artificially restricted the GA's search space and starved specific algorithms of fitness data. \\
\textbf{Resolution:} The \texttt{getEnemyTypesForFloor} function in \texttt{Archetypes.ts} was modified to unlock all eight enemy archetypes immediately from Floor 1. Furthermore, the spawn loop in \texttt{GameScreen.tsx} was overhauled to use a round-robin distribution strategy rather than randomized selection. This mathematically guarantees that at least one instance of every enemy archetype (and thus, every pathfinding algorithm) is spawned per map iteration, providing the GA with a universally diverse dataset.

\subsection{Analytics Accessibility During Pause}
\textbf{Issue:} The global pause menu's \texttt{z-index} blocked the React analytics dashboard overlay, preventing developers from inspecting telemetry during a frozen frame. \\
\textbf{Resolution:} The component hierarchy in \texttt{GameScreen.tsx} was adjusted so the analytics panel renders above the pause context layer. The backtick (\`) key listener was relocated to a globally scoped input handler that remains active while the simulation is paused.

\subsection{Correcting Metric Pollution}
\textbf{Issue:} The analytics snapshot loop was pulling data from dead, despawned, or nullified entities, polluting the metrics and displaying "ghost" cards in the UI. \\
\textbf{Resolution:} The snapshot timer in \texttt{GameScreen.tsx} now utilizes a strict \texttt{.filter((e) => e.isAlive)} map before dispatching data to the Zustand store. The live AI monitor cards in \texttt{AIAnalyticsPanel.tsx} were also updated to display the active algorithm using accurate, color-coded badges matching the \texttt{EnemyType}.

\section{Pathfinding Logic \& AI Deadlocks}

\subsection{Pathfinding State Mutation Consistency}
\textbf{Issue:} Upon entering the \texttt{ALERT} or \texttt{CHASING} state, enemies randomly sampled an algorithm from their genome weights, abandoning their archetype's default algorithm (e.g., a Ghost switching from DFS to A* mid-pursuit). \\
\textbf{Resolution:} The \texttt{getActiveAlgorithm} method in \texttt{EnemyBase.ts} was corrected. Enemies now strictly adhere to their assigned \texttt{ENEMY\_DEFAULT\_ALGORITHM}. The genome algorithm weights are now isolated strictly for evolutionary fitness tracking and mutation, and no longer govern runtime pathfinding overrides.

\subsection{Grid-Based Local Minima in DLS \& IDS}
\textbf{Issue:} The Ogre (Depth-Limited Search) and Bridge Heroine (Iterative Deepening Search) frequently became permanently stuck against walls. Because these algorithms utilize a generic \texttt{visited} state, "snake paths" generated by deep exploratory branches marked adjacent tiles as visited, mathematically blocking shorter, more direct paths. \\
\textbf{Resolution:} Both \texttt{DLS.ts} and \texttt{IDS.ts} were rewritten to utilize a \texttt{visitedDepth} map. Instead of a binary visited flag, the algorithm records the depth at which a node was discovered. If the search routine discovers a previously visited node via a shorter path (\texttt{nextDepth < prevDepth}), it is allowed to overwrite the history and continue exploring. This breaks the local minima deadlock while remaining mathematically identical to pure depth-first exploration.

\subsection{Hill Climbing Vantage Scoring (Distance Penalty)}
\textbf{Issue:} The Demon (Hill Climbing) refused to pursue the player, instead choosing to retreat into empty corners of the map. \\
\textbf{Resolution:} An analysis of \texttt{HillClimbing.ts} revealed a gradient inversion. The default scoring awarded $+3$ points per unit of distance away from walls. The distance penalty for being far from the player was only $-2$. Consequently, moving away from the player to avoid walls resulted in a net positive score. The heuristic gradient was severely steepened; moving closer to the target player now yields $+100$ base points minus a penalty multiplier of $20$, utterly overwhelming the wall-avoidance logic and forcing relentless pursuit until the optimal ranged distance is achieved.

\subsection{Asynchronous Path Following Stutter}
\textbf{Issue:} When alerted, enemies request paths asynchronously every 0.12 seconds. Because the path is generated relative to the time of request, enemies that moved forward one tile while waiting for the web-worker response would receive a path telling them to step backward onto the tile they just left. This drained their \texttt{moveAccumulator} budget to zero, causing them to vibrate in place. \\
\textbf{Resolution:} The path reception logic in \texttt{EnemyBase.ts} was patched. Upon receiving a resolved path, the enemy now iterates through the path array to dynamically discover its precise, current topological location. It then sets its \texttt{pathIndex} to \texttt{startIndex + 1}, ensuring it never wastes movement cycles re-navigating to an occupied tile.

\subsection{Map Selection Persistence}
\textbf{Issue:} Despite explicit user selection, the engine would transition from the chosen Tiled JSON map to a procedurally generated BSP dungeon upon reaching Floor 3 (Iteration 2). \\
\textbf{Resolution:} The procedural switch within \texttt{DungeonGenerator.ts} was governed by a hardcoded \texttt{if (floor === 1 || floor === 2)} limit. This condition was completely removed. The generator now unconditionally prioritizes the \texttt{selectedMap} variable regardless of the current evolutionary iteration.

\section{Conclusion}
The completion of Phase 6 definitively stabilizes the EvoCracker environment. By resolving fundamental mathematical oversights in the search algorithms and enforcing strict separation between visual telemetry and underlying genetic data, the simulation now provides a pristine, deterministically accurate dataset for continuous evolutionary learning.

\end{document}
